
\documentclass[10pt,aspectratio=169]{beamer}

\usetheme{Madrid}
\usecolortheme{default}

\title{Discussion}
\subtitle{School Spending and Students' Achievement}
\date{}

\begin{document}

% Slide 1: Questions
\begin{frame}{Discussion -- Questions}
\begin{enumerate}
    \item \textbf{Budget channel:}
You argue that the 2014-2016 oil bust must affects municipalities uniquely through a budget channel since Macaé is excluded. But it seems plausible that such a downturn in the Campos Basin could also affects students through household incomes or migration. How do you plan to disentangle those effects ?

    \vspace{0.3cm}

    \item \textbf{Brazilian recession:}
    How can you isolate the oil shock from the broader
    economic recession in Brazil during 2014--2016? (question Eve à compléter)

    \vspace{0.3cm}

    \item \textbf{Question 3:}
   How credible is the parallel trends assumption, given the structural differences between oil-dependent and other municipalities? Would the results be robust to using a more comparable control group?

    \vspace{0.3cm}

    \item \textbf{Question 4:}
    To be completed.
\end{enumerate}
\end{frame}

% Slide 2: Feedback
\begin{frame}{Discussion -- Feedback}
\begin{itemize}
    

    \vspace{0.3cm}

    \item \textbf{Identification:}
    

    \vspace{0.3cm}

    \item \textbf{Interpretation:}
    
\end{itemize}
\end{frame}

\end{document}
